\subsection*{15.7　初等奇点的几何变式}
\begingroup
\def\OO{\mathcal O}\def\m{\mathfrak m}\def\mm{\mathfrak m}\def\Sym{\operatorname{Sym}}
应当提出15.6的“几何”变式。特别地，将得到如下结论：设$X$是域$k$上局部有限型的预概形，$x\in X$，并令$d=\dim_x X$。$x$是$X$的重数$n$的\emph{几何}初等奇点，当且仅当$x$有一个开邻域$U$，可以把$U$嵌为$k$上连通的、维数为$d+1$的光滑预概形$X'$的子预概形，且$U$由一个方程$\varphi=0$定义，$x$是$\varphi$的重数$n$的几何普通奇异零点。\footnote{法文用$d=\dim_x X$、$d+1$和底域$k$，英译首句误把$d$作$n$，并将后一处底域作$K$。法文在判别的左边只写“初等奇点”，而右边“普通奇异零点”上方手写添了“几何”。正文按标题及右边补明左边的“几何”；绝对意义不可直接代入，见后附反例。}

此外，我们在15.5中本应已经说明：这就是说，$\varphi$的主部$d^n\varphi$在$x$处的值（先验地是$\mathcal P^n_{X'/k}\otimes\kappa(x)$的一个元素，更确切地说是其增广理想$I$中的一个元素）实际上落在
\[
I^n=\Sym^n_{\kappa(x)}
\bigl(\Omega^1_{X'/k}\otimes\kappa(x)\bigr)
\]
中，并且作为这样的元素，它是一个光滑型。\footnote{法文可辨主部符号，原对称幂标作$\Sym^n_{\OO_U}$；此处把纤维上的对称幂明确写在$\kappa(x)$上。截断主部环中$I^{n+1}=0$，故最后一层$I^n/I^{n+1}$可记为$I^n$。按重数恰为$n$，其次数$n$的型须非零；这是译校说明，不是再加一条原文译句。}
〔注意：15.1中漏说了，型的光滑点集是开集。〕

还可注意到，这样的点在非光滑点集中是孤立的；它几何正规（分别地，几何约化）当且仅当$d\ge2$（分别地，$d\ge1$）。若$d=0$且$\kappa(x)=k$（即$x$是孤立的$k$-有理点），则$x$为重数$n$的初等奇点，意即$\OO_{X,x}$作为$k$-代数同构于$k[T]/(T^n)$。

\newpage
若$d=1$且$x$在$k$上有理，那么重数$n$的初等奇点，在经典术语中就是“有$n$条相异切线的普通奇点”。（我们在15.2中本应已经讲明：若$x$是$X$的一个$k$-有理点，则重数$n$的初等奇点在绝对意义和相对意义下是同一概念；若$\kappa(x)$是$k$的有限可分扩张，这也仍然成立。作者还猜测“有限”大概可省。\footnote{作者的猜测在允许非闭点并采用15.2的有理提升定义时不成立；见后文。}此事值得列为一条推论或命题。）

\paragraph{作者的编排附记。}
注意：我在第20节中曾列入一小节，论述沿闭子预概形$Y$吹起预概形$X$，其中增广$\OO_X\to\OO_Y$是重数$n$的初等增广（15.6已提到这项结果）。这里值得一提的是，这一小节简短而切题；它只用到第17节的一般光滑性结果和17.15.2开头的定义（甚至后者也可避免）。完全可以现在就把这些结果纳入，例如放到上文的定义15.1之后；这样便可以把15.2中的一组“几何”结果以及接着的15.3至15.7与这些结果分开，另编为17.16一节。

不论怎样，我写着写着发现，原拟论述光滑型的新第15小节，实际上至少要拆成两节，二者相当不同且彼此独立：一节汇集任意$n$时关于“重数$n$的初等”对象的一般结论；另一节刻画光滑的二次型，而从前一节只取用15.1，也就是实际上只取用型之光滑性的定义。

\paragraph{译校说明（非原文）。}
\paragraph{截断主部不是除以$n!$的导数。}
设$X'=\Spec R$，$J=\ker(R\otimes_kR\to R)$，$K=\kappa(x)$。用左因子在$x$求值，对$(R\otimes_k R)/J^{n+1}$基变换，得到$(K\otimes_k R)/\m_z^{n+1}$，其中$z$是$X'_K$的规范$K$-有理提升，$\m_z$为其点理想。该商也等于$\OO_{X'_K,z}/\m_z^{n+1}$，且$1\otimes\varphi$给出所论主部。落在$I^n$的含义是所有低于$n$阶的主部消失，由光滑环境的正则性，次数$n$的分量由切余空间的$n$次对称幂给出。它非零且切出光滑的射影超曲面，恰是有理提升处“普通零点”的这里所需条件。在任何特征均不需要以$n!$作除法。参见 \href{https://stacks.math.columbia.edu/tag/09CH}{主部代数}、\href{https://stacks.math.columbia.edu/tag/08S2}{对角余法模}和\href{https://stacks.math.columbia.edu/tag/00NO}{正则局部环的分次环}。

\paragraph{几何孤立性与局部超曲面。}
取代数闭域$\Omega\supseteq\kappa(x)$，在$X_\Omega$中取$\Omega$-有理提升$z$。若$x$是几何初等奇点，则由15.6的局部结论$\OO_{X_\Omega,z}$的穿孔谱正则；因$\Omega$完美，$X_\Omega$在这些点也光滑。因而$X_\Omega$的非光滑闭集在有理闭点$z$处局部维数为零。非光滑集按\href{https://stacks.math.columbia.edu/tag/02V4}{光滑轨迹的基变性质}来自$X$的非光滑集。利用\href{https://stacks.math.columbia.edu/tag/00P4}{概形在一点处的维数在域扩张下不变}，得到该闭集在$x$处维数为零，因此$x$是其中的孤立闭点。这里的域扩张维数公式使用的是\emph{概形在点处的维数}，并非任意非闭点的局部环维数公式。几何正规和几何约化的阈值可对代数闭域上的有理提升用15.6检验，再按域扩张下降。

现取仿射嵌入$X\hookrightarrow\mathbf A^M_k$，$K=\kappa(x)$。前段表明此时$x$为闭点，几何切余空间维数为$d+1$。从$X$的定义理想中选取$M-d-1$个方程，使其微分的值构成环境至$X$的切余空间映射的核的一组基。由相对 Jacobian 判据，这些方程切出一个在$x$附近光滑、维数为$d+1$的环境；缩小到含$x$的连通分支。环境中$X$的余下定义理想，扩张至规范$K$-有理提升的局部环后，由15.6的初始理想论证为主理想。环境局部环至该提升的局部环忠实平坦，两者剩余域均为$K$，故余理想的最小生成元数不变；由 Nakayama 引理，它在$x$处主生成。缩小后即得所求方程；逆向由主部计算与15.6得到。
\paragraph{域扩张的限定与反例。}
\paragraph{有限可分与零维有理情形。}
若$\kappa(x)/k$有限可分，基变至$\kappa(x)$是\href{https://stacks.math.columbia.edu/tag/02GH}{étale 基变换}。在规范有理提升$z$处，局部同态平坦，且因$\kappa(x)\otimes_k\kappa(x)$在对角处局部化后等于$\kappa(x)$，极大理想由原极大理想延拓生成。故极大理想分次环等同，绝对与几何条件等价；$k$-有理点情形尤然。若$k$-有理且$d=0$，由15.6，极大理想有一个幂零生成元$\pi$，消失指数为$n$。已有结构映射$k\to\OO_{X,x}$和剩余域同构，故逐次按极大理想幂展开，所得映射$k[T]/(T^n)\to\OO_{X,x}$（$T\mapsto\pi$）满射，两侧均为$n$维$k$-向量空间，于是同构。$d=1$有理时，光滑二元次数$n$型等价于在代数闭包中分解成$n$个互异的线性因子；“相异切线”不要求每条都在$k$上定义。

\paragraph{不能删掉“有限”。}
令$X=\Spec k[t,u,v]/(uv)$，$x=(u,v)$为奇异直线的泛点。$\kappa(x)=k(t)$是可分生成扩张，绝对分次环为$k(t)[U,V]/(UV)$；射影切锥是两个光滑点，所以$x$绝对初等二次。取$K=k(a)$，其中$a$超越，令$t$映到$a$，则在$X_K$中由该嵌入确定的有理提升是$z=(t-a,u,v)$。它的分次环为$K[T,U,V]/(UV)$，射影切锥为两条在$[1:0:0]$相交的直线，不光滑。因此$x$不是几何初等二次奇点。

\paragraph{纯不可分时，也不能反过来以绝对替代几何。}
设$\operatorname{char}k=p>0$，$a\notin k^p$，取$X=\Spec(k[t]/(t^p-a))$。它的局部环是域，绝对正则。扩域至含$a^{1/p}$的$K$后，有理提升的局部环为$K[T]/(T^p)$；其次数$p$初始型只在$\mathbf P^0_K$截出空集，故按15.2是重数$p$的几何初等奇点。此例同时说明正文首句为何须写“几何”。
\endgroup
